% Wolf Prize in Mathematics — Laureates (1978–2024)
% Compact tabular layout mirroring fields_medal_laureates_beamer.tex
% Cross-award badges per math_awards_allinone.tex
% Compile with: xelatex wolf_prize_laureates_beamer.tex

\documentclass[aspectratio=169,9pt]{beamer}

% ---- Theme ----
\usetheme{default}
\usecolortheme{default}
\setbeamertemplate{navigation symbols}{}
\setbeamertemplate{footline}{%
  \hfill{\tiny\color{mutedgray}\insertframenumber/\inserttotalframenumber}\hspace*{6pt}\vspace{4pt}%
}
\setbeamertemplate{frametitle}{%
  \vspace{4pt}%
  \noindent\hspace*{8pt}{\large\bfseries\color{wolfclr}\insertframetitle}%
  \hfill
  \raisebox{-1pt}{\footnotesize\color{mutedgray} 沃尔夫数学奖 Wolf Prize}%
  \hspace{8pt}%
  \vspace{2pt}%
  \par\noindent\hspace*{8pt}\textcolor{wolfbg}{\rule{0.95\textwidth}{0.6pt}}%
  \vspace{-2pt}%
}

% ---- Fonts & Language ----
\usepackage{fontspec}
\usepackage{xeCJK}
\setCJKmainfont{PingFang SC}[BoldFont=PingFang SC Semibold]
\setCJKsansfont{PingFang SC}
\setmainfont{Helvetica Neue}[BoldFont=Helvetica Neue Bold]
\setsansfont{Helvetica Neue}

% ---- Packages ----
\usepackage{booktabs}
\usepackage{array}
\usepackage{tikz}
\usepackage{amssymb}
\usepackage{colortbl}
\usepackage{graphicx}

% ---- Colors ----
\definecolor{headerblue}{RGB}{41,65,122}
\definecolor{yearcolor}{RGB}{140,90,0}
\definecolor{fieldcolor}{RGB}{30,100,60}
\definecolor{mutedgray}{RGB}{100,100,100}
\definecolor{slidebg}{RGB}{255,253,248}
% ---- Per-award brand colors (match math_awards_allinone) ----
\definecolor{fieldsclr}{RGB}{41,65,122}   % Fields — 深蓝
\definecolor{wolfclr}{RGB}{150,100,0}     % Wolf   — 金棕（主色）
\definecolor{abelclr}{RGB}{150,40,90}     % Abel   — 紫红
\definecolor{chernclr}{RGB}{15,110,92}    % Chern  — 翡翠绿
\definecolor{wolfbg}{RGB}{240,228,200}    % 表头底色

% ---- Beamer Colors ----
\setbeamercolor{background canvas}{bg=slidebg}
\setbeamercolor{title}{fg=wolfclr}

% ---- Custom Commands ----
% 奖项徽标：追加在姓名后，标注该得主还获得的其他顶级奖项
\newcommand{\fieldsbadge}{\textsuperscript{\normalfont\tiny\color{fieldsclr}$\blacklozenge$\kern-0.5pt F}}
\newcommand{\abelbadge}{\textsuperscript{\normalfont\tiny\color{abelclr}$\blacktriangle$\kern-0.5pt A}}
\newcommand{\chernbadge}{\textsuperscript{\normalfont\tiny\color{chernclr}$\blacksquare$\kern-0.5pt C}}

% Table row for a laureate
% #1=Name(+badges), #2=Year, #3=Birth-Death, #4=Nationality, #5=Institution, #6=Field
% Name cell carries a \cn{中文名} suffix rendered on its own line in muted gray.
\newcommand{\cn}[1]{\newline{\scriptsize\color{mutedgray}#1}}
\newcommand{\lrow}[6]{%
  {\bfseries\color{wolfclr}#1} & {\color{yearcolor}#2} & {#3} & {#4} & {\scriptsize #5} & {\scriptsize\color{fieldcolor}#6} \\[2.5pt]
}

% Year section header inside table
\newcommand{\yearrow}[1]{%
  \rowcolor{wolfbg}
  \multicolumn{6}{l}{\raisebox{0pt}[1.1em][0.3em]{\bfseries\color{yearcolor}\normalsize #1}} \\[0.5pt]
}

% ---- Document ----
\begin{document}

% ============================================================
% Title Slide
% ============================================================
\begin{frame}
\vfill
\begin{center}
  {\Huge\bfseries\color{wolfclr} 沃尔夫数学奖获奖者名录}\\[8pt]
  {\Large\color{wolfclr} Wolf Prize in Mathematics}\\[6pt]
  {\normalsize\color{mutedgray} 1978–2024 · 以色列沃尔夫基金会 · 无年龄限制}\\[12pt]
  {\small\color{mutedgray} 按颁奖年份排列}\\[8pt]
  {\footnotesize 交叉获奖标注：%
    {\color{fieldsclr}$\blacklozenge$F} 菲尔兹奖\quad
    {\color{abelclr}$\blacktriangle$A} 阿贝尔奖\quad
    {\color{chernclr}$\blacksquare$C} 陈省身奖章}
\end{center}
\vspace{10pt}
\begin{center}
  {\small\color{mutedgray}\itshape 「表彰对数学终身的杰出贡献」—— 与菲尔兹奖并称，无年龄限制，每年评选}
\end{center}
\vfill
\end{frame}

% Table slide 1: 1978–1980
\begin{frame}{1978–1980}
\vspace{-2pt}
\centering
\scriptsize
\begin{tabular}{@{}
  >{\raggedright}m{3.7cm}
  >{\centering}m{0.95cm}
  >{\centering}m{1.05cm}
  >{\raggedright}m{1.35cm}
  >{\raggedright}m{2.35cm}
  >{\raggedright\arraybackslash}m{2.55cm}
@{}}
\toprule
\rowcolor{wolfbg}
\textbf{姓名} & \textbf{年份} & \textbf{生卒} & \textbf{国籍} & \textbf{机构} & \textbf{研究领域} \\
\midrule
\lrow{Israel Gelfand\cn{以色列·盖尔范德}}{1978}{1913–2009}{苏联}{Moscow State University}{泛函分析、群表示论、积分几何；首届得主}
\lrow{Carl Ludwig Siegel\cn{卡尔·西格尔}}{1978}{1896–1981}{西德}{University of Göttingen}{数论、多复变函数、天体力学}
\lrow{Jean Leray\cn{让·勒雷}}{1979}{1906–1998}{法国}{Collège de France}{层论、谱序列、拓扑度理论}
\lrow{André Weil\cn{安德烈·韦伊}}{1979}{1906–1998}{法国}{Institute for Advanced Study}{代数几何基础、数论；Bourbaki 核心}
\lrow{Henri Cartan\cn{亨利·嘉当}}{1980}{1904–2008}{法国}{Université Paris-Sud}{同调代数、层论、多复变；Bourbaki 创始}
\lrow{Andrey Kolmogorov\cn{安德烈·柯尔莫哥洛夫}}{1980}{1903–1987}{苏联}{Moscow State University}{概率论公理化、湍流、动力系统}
\bottomrule
\end{tabular}
\end{frame}

% Table slide 2: 1981–1983/84
\begin{frame}{1981–1983/84}
\vspace{-2pt}
\centering
\scriptsize
\begin{tabular}{@{}
  >{\raggedright}m{3.7cm}
  >{\centering}m{0.95cm}
  >{\centering}m{1.05cm}
  >{\raggedright}m{1.35cm}
  >{\raggedright}m{2.35cm}
  >{\raggedright\arraybackslash}m{2.55cm}
@{}}
\toprule
\rowcolor{wolfbg}
\textbf{姓名} & \textbf{年份} & \textbf{生卒} & \textbf{国籍} & \textbf{机构} & \textbf{研究领域} \\
\midrule
\lrow{Lars Ahlfors\fieldsbadge\cn{拉尔斯·阿尔福斯}}{1981}{1907–1996}{芬兰/美国}{Harvard University}{黎曼曲面、拟共形映射；亦获 F1936}
\lrow{Oscar Zariski\cn{奥斯卡·扎里斯基}}{1981}{1899–1986}{美国}{Harvard University}{代数几何基础、Zariski 拓扑}
\lrow{Hassler Whitney\cn{哈斯勒·惠特尼}}{1982}{1907–1989}{美国}{Institute for Advanced Study}{微分拓扑、奇点理论、拟阵理论}
\lrow{Mark Krein\cn{马克·克列因}}{1982}{1907–1989}{苏联}{Odessa University}{泛函分析、算子理论、矩问题}
\lrow{Shiing-Shen Chern\cn{陈省身}}{1983/84}{1911–2004}{中国/美国}{University of California, Berkeley}{微分几何、陈类、Chern-Simons；奖章即以其命名}
\lrow{Paul Erdős\cn{保罗·埃尔德什}}{1983/84}{1913–1996}{匈牙利}{Hungarian Academy of Sciences}{组合数学、数论、图论}
\bottomrule
\end{tabular}
\end{frame}

% Table slide 3: 1984/85–1987
\begin{frame}{1984/85–1987}
\vspace{-2pt}
\centering
\scriptsize
\begin{tabular}{@{}
  >{\raggedright}m{3.7cm}
  >{\centering}m{0.95cm}
  >{\centering}m{1.05cm}
  >{\raggedright}m{1.35cm}
  >{\raggedright}m{2.35cm}
  >{\raggedright\arraybackslash}m{2.55cm}
@{}}
\toprule
\rowcolor{wolfbg}
\textbf{姓名} & \textbf{年份} & \textbf{生卒} & \textbf{国籍} & \textbf{机构} & \textbf{研究领域} \\
\midrule
\lrow{Kunihiko Kodaira\fieldsbadge\cn{小平邦彦}}{1984/85}{1915–1997}{日本}{University of Tokyo / Gakushuin University}{复流形、Hodge 理论；亦获 F1954}
\lrow{Hans Lewy\cn{汉斯·列维}}{1984/85}{1904–1988}{美国}{UC Berkeley}{偏微分方程}
\lrow{Samuel Eilenberg\cn{塞缪尔·艾伦伯格}}{1986}{1913–1998}{波兰/美国}{Columbia University}{同调代数、范畴论}
\lrow{Atle Selberg\fieldsbadge\cn{阿特勒·塞尔伯格}}{1986}{1917–2007}{挪威/美国}{Institute for Advanced Study}{筛法、Selberg 迹公式；亦获 F1950}
\lrow{Kiyoshi Itō\cn{伊藤清}}{1987}{1915–2008}{日本}{Kyoto University}{随机微积分、Itō 引理；首届高斯奖}
\lrow{Peter Lax\abelbadge\cn{彼得·拉克斯}}{1987}{1926–}{匈牙利/美国}{Courant Institute of Mathematical Sciences, New York University}{PDE、Lax 对；亦获 A2005}
\bottomrule
\end{tabular}
\end{frame}

% Table slide 4: 1988–1990
\begin{frame}{1988–1990}
\vspace{-2pt}
\centering
\scriptsize
\begin{tabular}{@{}
  >{\raggedright}m{3.7cm}
  >{\centering}m{0.95cm}
  >{\centering}m{1.05cm}
  >{\raggedright}m{1.35cm}
  >{\raggedright}m{2.35cm}
  >{\raggedright\arraybackslash}m{2.55cm}
@{}}
\toprule
\rowcolor{wolfbg}
\textbf{姓名} & \textbf{年份} & \textbf{生卒} & \textbf{国籍} & \textbf{机构} & \textbf{研究领域} \\
\midrule
\lrow{Friedrich Hirzebruch\cn{弗里德里希·希策布鲁赫}}{1988}{1927–2012}{西德}{University of Bonn / Max Planck Institute}{拓扑、代数几何、HRR 定理；创立 MPIM}
\lrow{Lars Hörmander\fieldsbadge\cn{拉尔斯·赫尔曼德}}{1988}{1931–2012}{瑞典}{Lund University}{线性偏微分算子；亦获 F1962}
\lrow{Alberto Calderón\cn{阿尔韦托·卡尔德龙}}{1989}{1920–1998}{阿根廷/美国}{University of Chicago}{奇异积分、Calderón-Zygmund 理论}
\lrow{John Milnor\fieldsbadge\abelbadge\cn{约翰·米尔诺}}{1989}{1931–}{美国}{SUNY Stony Brook}{微分拓扑、异国球面；亦获 F1962 / A2011}
\lrow{Ennio De Giorgi\cn{恩尼奥·德乔尔吉}}{1990}{1928–1996}{意大利}{Scuola Normale Superiore di Pisa}{极小曲面、变分法、Hilbert 第 19 问题}
\lrow{Ilya Piatetski-Shapiro\cn{伊利亚·皮亚捷茨基-夏皮罗}}{1990}{1929–2009}{苏联/以色列}{Tel Aviv University / Yale}{自守形式、L 函数、表示论}
\bottomrule
\end{tabular}
\end{frame}

% Table slide 5: 1992–1994/95
\begin{frame}{1992–1994/95}
\vspace{-2pt}
\centering
\scriptsize
\begin{tabular}{@{}
  >{\raggedright}m{3.7cm}
  >{\centering}m{0.95cm}
  >{\centering}m{1.05cm}
  >{\raggedright}m{1.35cm}
  >{\raggedright}m{2.35cm}
  >{\raggedright\arraybackslash}m{2.55cm}
@{}}
\toprule
\rowcolor{wolfbg}
\textbf{姓名} & \textbf{年份} & \textbf{生卒} & \textbf{国籍} & \textbf{机构} & \textbf{研究领域} \\
\midrule
\lrow{Lennart Carleson\abelbadge\cn{伦纳特·卡尔松}}{1992}{1928–}{瑞典}{Uppsala University / UCLA}{调和分析、Carleson 定理；亦获 A2006}
\lrow{John G. Thompson\fieldsbadge\abelbadge\cn{约翰·汤普森}}{1992}{1932–}{美国}{University of Cambridge}{有限群论、奇阶定理；亦获 F1970 / A2008}
\lrow{Mikhail Gromov\abelbadge\cn{米哈伊尔·格罗莫夫}}{1993}{1943–}{俄罗斯/法国}{Institut des Hautes Études Scientifiques (IHÉS)}{黎曼几何、辛几何、几何群论；亦获 A2009}
\lrow{Jacques Tits\abelbadge\cn{雅克·蒂茨}}{1993}{1930–2021}{比利时/法国}{Collège de France}{群论、建筑理论；亦获 A2008}
\lrow{Jürgen Moser\cn{于尔根·莫泽}}{1994/95}{1928–1999}{德国/瑞士}{ETH Zürich}{KAM 理论、动力系统}
\bottomrule
\end{tabular}
\end{frame}

% Table slide 6: 1995/96–1999
\begin{frame}{1995/96–1999}
\vspace{-2pt}
\centering
\scriptsize
\begin{tabular}{@{}
  >{\raggedright}m{3.7cm}
  >{\centering}m{0.95cm}
  >{\centering}m{1.05cm}
  >{\raggedright}m{1.35cm}
  >{\raggedright}m{2.35cm}
  >{\raggedright\arraybackslash}m{2.55cm}
@{}}
\toprule
\rowcolor{wolfbg}
\textbf{姓名} & \textbf{年份} & \textbf{生卒} & \textbf{国籍} & \textbf{机构} & \textbf{研究领域} \\
\midrule
\lrow{Robert Langlands\abelbadge\cn{罗伯特·朗兰兹}}{1995/96}{1936–}{加拿大/美国}{Institute for Advanced Study}{Langlands 纲领；亦获 A2018}
\lrow{Andrew Wiles\abelbadge\cn{安德鲁·怀尔斯}}{1995/96}{1953–}{英国}{Princeton University}{费马大定理；亦获 A2016}
\lrow{Joseph B. Keller\cn{约瑟夫·凯勒}}{1996/97}{1923–2016}{美国}{Stanford University}{应用数学、渐近方法、几何衍射理论}
\lrow{Yakov Sinai\abelbadge\cn{雅科夫·西奈}}{1996/97}{1935–}{俄罗斯/美国}{Landau Institute / Princeton}{动力系统、遍历理论；亦获 A2014}
\lrow{László Lovász\abelbadge\cn{拉斯洛·洛瓦兹}}{1999}{1948–}{匈牙利}{Eötvös Loránd University / Yale}{组合数学、图论、LLL 算法；亦获 A2021}
\lrow{Elias M. Stein\cn{埃利亚斯·施泰因}}{1999}{1931–2018}{美国}{Princeton University}{调和分析、奇异积分；培养多位菲尔兹奖得主}
\bottomrule
\end{tabular}
\end{frame}

% Table slide 7: 2000–2002/03
\begin{frame}{2000–2002/03}
\vspace{-2pt}
\centering
\scriptsize
\begin{tabular}{@{}
  >{\raggedright}m{3.7cm}
  >{\centering}m{0.95cm}
  >{\centering}m{1.05cm}
  >{\raggedright}m{1.35cm}
  >{\raggedright}m{2.35cm}
  >{\raggedright\arraybackslash}m{2.55cm}
@{}}
\toprule
\rowcolor{wolfbg}
\textbf{姓名} & \textbf{年份} & \textbf{生卒} & \textbf{国籍} & \textbf{机构} & \textbf{研究领域} \\
\midrule
\lrow{Raoul Bott\cn{拉乌尔·博特}}{2000}{1923–2005}{匈牙利/美国}{Harvard University}{拓扑、微分几何、Bott 周期性定理}
\lrow{Jean-Pierre Serre\fieldsbadge\abelbadge\cn{让-皮埃尔·塞尔}}{2000}{1926–}{法国}{Collège de France}{代数拓扑、代数几何、数论；亦获 F1954 / A2003}
\lrow{Vladimir Arnold\cn{弗拉基米尔·阿诺尔德}}{2001}{1937–2010}{俄罗斯}{Steklov Mathematical Institute / Paris-Dauphine}{动力系统、奇点理论、KAM 理论}
\lrow{Saharon Shelah\cn{萨哈龙·谢拉赫}}{2001}{1945–}{以色列}{Hebrew University of Jerusalem}{数理逻辑、集合论、模型论}
\lrow{Mikio Sato\cn{佐藤幹夫}}{2002/03}{1928–2023}{日本}{Kyoto University}{代数分析、D 模理论、Sato 超函数}
\lrow{John Tate\abelbadge\cn{约翰·泰特}}{2002/03}{1925–2019}{美国}{University of Texas at Austin}{数论、Tate 论题；亦获 A2010}
\bottomrule
\end{tabular}
\end{frame}

% Table slide 8: 2005–2006/07
\begin{frame}{2005–2006/07}
\vspace{-2pt}
\centering
\scriptsize
\begin{tabular}{@{}
  >{\raggedright}m{3.7cm}
  >{\centering}m{0.95cm}
  >{\centering}m{1.05cm}
  >{\raggedright}m{1.35cm}
  >{\raggedright}m{2.35cm}
  >{\raggedright\arraybackslash}m{2.55cm}
@{}}
\toprule
\rowcolor{wolfbg}
\textbf{姓名} & \textbf{年份} & \textbf{生卒} & \textbf{国籍} & \textbf{机构} & \textbf{研究领域} \\
\midrule
\lrow{Grigory Margulis\fieldsbadge\abelbadge\cn{格里戈里·马尔古利斯}}{2005}{1946–}{俄罗斯/美国}{Yale University}{Lie 群、格、刚性理论；亦获 F1978 / A2020}
\lrow{Sergei Novikov\fieldsbadge\cn{谢尔盖·诺维科夫}}{2005}{1938–2024}{俄罗斯}{University of Maryland / Landau Institute}{拓扑、Pontryagin 类、孤立子理论；亦获 F1970}
\lrow{Stephen Smale\fieldsbadge\cn{斯蒂芬·斯梅尔}}{2006/07}{1930–}{美国}{UC Berkeley / Toyota Tech / City Univ. of HK}{拓扑、动力系统、数值分析；亦获 F1966}
\lrow{Hillel Furstenberg\abelbadge\cn{希勒尔·富斯滕伯格}}{2006/07}{1935–}{以色列/美国}{Hebrew University of Jerusalem}{遍历理论、概率、组合数论；亦获 A2020}
\bottomrule
\end{tabular}
\end{frame}

% Table slide 9: 2008–2010
\begin{frame}{2008–2010}
\vspace{-2pt}
\centering
\scriptsize
\begin{tabular}{@{}
  >{\raggedright}m{3.7cm}
  >{\centering}m{0.95cm}
  >{\centering}m{1.05cm}
  >{\raggedright}m{1.35cm}
  >{\raggedright}m{2.35cm}
  >{\raggedright\arraybackslash}m{2.55cm}
@{}}
\toprule
\rowcolor{wolfbg}
\textbf{姓名} & \textbf{年份} & \textbf{生卒} & \textbf{国籍} & \textbf{机构} & \textbf{研究领域} \\
\midrule
\lrow{Pierre Deligne\fieldsbadge\abelbadge\cn{皮埃尔·德利涅}}{2008}{1944–}{比利时}{Institute for Advanced Study}{代数几何、Weil 猜想；亦获 F1978 / A2013}
\lrow{Phillip Griffiths\chernbadge\cn{菲利普·格里菲斯}}{2008}{1938–}{美国}{Institute for Advanced Study}{代数几何、Hodge 理论；亦获 C2014}
\lrow{David Mumford\fieldsbadge\cn{大卫·芒福德}}{2008}{1937–}{美国}{Brown University / Harvard}{代数几何、模空间；亦获 F1974}
\lrow{Shing-Tung Yau\fieldsbadge\cn{丘成桐}}{2010}{1949–}{中国/美国}{Harvard University}{Calabi 猜想、几何分析；亦获 F1982；首位华人双奖}
\lrow{Dennis Sullivan\abelbadge\cn{丹尼斯·苏利文}}{2010}{1941–}{美国}{SUNY Stony Brook / CUNY}{拓扑、有理同伦论、复动力系统；亦获 A2022}
\bottomrule
\end{tabular}
\end{frame}

% Table slide 10: 2012–2015
\begin{frame}{2012–2015}
\vspace{-2pt}
\centering
\scriptsize
\begin{tabular}{@{}
  >{\raggedright}m{3.7cm}
  >{\centering}m{0.95cm}
  >{\centering}m{1.05cm}
  >{\raggedright}m{1.35cm}
  >{\raggedright}m{2.35cm}
  >{\raggedright\arraybackslash}m{2.55cm}
@{}}
\toprule
\rowcolor{wolfbg}
\textbf{姓名} & \textbf{年份} & \textbf{生卒} & \textbf{国籍} & \textbf{机构} & \textbf{研究领域} \\
\midrule
\lrow{Michael Aschbacher\cn{迈克尔·阿施巴赫}}{2012}{1944–}{美国}{Caltech}{有限群论、有限单群分类}
\lrow{Luis Caffarelli\abelbadge\cn{路易斯·卡法雷利}}{2012}{1948–}{阿根廷/美国}{University of Texas at Austin}{非线性 PDE、自由边界问题；亦获 A2023}
\lrow{George Mostow\cn{乔治·莫斯托}}{2013}{1923–2017}{美国}{Yale University}{李群、Mostow 刚性定理}
\lrow{Michael Artin\cn{迈克尔·阿廷}}{2013}{1934–}{美国}{MIT}{代数几何、交换代数、非交换环}
\lrow{Peter Sarnak\cn{彼得·萨纳克}}{2014}{1953–}{南非/美国}{Princeton University / IAS}{数论、自守形式、随机矩阵理论}
\lrow{James Arthur\cn{詹姆斯·阿瑟}}{2015}{1944–}{加拿大}{University of Toronto}{自守形式表示论、Arthur-Selberg 迹公式}
\bottomrule
\end{tabular}
\end{frame}

% Table slide 11: 2017–2019
\begin{frame}{2017–2019}
\vspace{-2pt}
\centering
\scriptsize
\begin{tabular}{@{}
  >{\raggedright}m{3.7cm}
  >{\centering}m{0.95cm}
  >{\centering}m{1.05cm}
  >{\raggedright}m{1.35cm}
  >{\raggedright}m{2.35cm}
  >{\raggedright\arraybackslash}m{2.55cm}
@{}}
\toprule
\rowcolor{wolfbg}
\textbf{姓名} & \textbf{年份} & \textbf{生卒} & \textbf{国籍} & \textbf{机构} & \textbf{研究领域} \\
\midrule
\lrow{Charles Fefferman\fieldsbadge\cn{查尔斯·费弗曼}}{2017}{1949–}{美国}{Princeton University}{调和分析、多复变、PDE；亦获 F1978}
\lrow{Richard Schoen\cn{理查德·舍恩}}{2017}{1950–}{美国}{UC Irvine / Stanford}{微分几何、几何分析、Yamabe 问题}
\lrow{Alexander Beilinson\cn{亚历山大·贝林森}}{2018}{1957–}{俄罗斯/美国}{University of Chicago}{代数几何、表示论、Beilinson 猜想}
\lrow{Vladimir Drinfeld\fieldsbadge\cn{弗拉基米尔·德林费尔德}}{2018}{1954–}{乌克兰/美国}{University of Chicago}{Langlands 纲领、量子群；亦获 F1990}
\lrow{Jean-François Le Gall\cn{让-弗朗索瓦·勒加尔}}{2019}{1959–}{法国}{Université Paris-Sud / Institut Universitaire de France}{随机过程、布朗运动、随机几何}
\lrow{Gregory Lawler\cn{格雷戈里·劳勒}}{2019}{1955–}{美国}{University of Chicago}{随机过程、SLE 理论}
\bottomrule
\end{tabular}
\end{frame}

% Table slide 12: 2020–2024
\begin{frame}{2020–2024}
\vspace{-2pt}
\centering
\scriptsize
\begin{tabular}{@{}
  >{\raggedright}m{3.7cm}
  >{\centering}m{0.95cm}
  >{\centering}m{1.05cm}
  >{\raggedright}m{1.35cm}
  >{\raggedright}m{2.35cm}
  >{\raggedright\arraybackslash}m{2.55cm}
@{}}
\toprule
\rowcolor{wolfbg}
\textbf{姓名} & \textbf{年份} & \textbf{生卒} & \textbf{国籍} & \textbf{机构} & \textbf{研究领域} \\
\midrule
\lrow{Simon Donaldson\fieldsbadge\cn{西蒙·唐纳森}}{2020}{1957–}{英国}{Imperial College London / SUNY Stony Brook}{四维流形拓扑、规范理论；亦获 F1986}
\lrow{Yakov Eliashberg\cn{雅科夫·埃利亚什贝格}}{2020}{1946–}{俄罗斯/美国}{Stanford University}{辛几何、接触几何、h 原理}
\lrow{George Lusztig\cn{乔治·卢斯蒂格}}{2022}{1946–}{罗马尼亚/美国}{MIT}{表示论、代数群、Hecke 代数、量子群}
\lrow{Ingrid Daubechies\cn{英格丽·多贝西}}{2023}{1954–}{比利时/美国}{Duke University}{小波理论、图像压缩；首位女性沃尔夫数学奖得主}
\lrow{Noga Alon\cn{诺加·阿隆}}{2024}{1956–}{以色列}{Tel Aviv University}{组合数学、图论、组合 Nullstellensatz}
\lrow{Adi Shamir\cn{阿迪·萨米尔}}{2024}{1952–}{以色列}{Weizmann Institute of Science}{RSA 加密算法、密码学；亦获 2002 图灵奖}
\bottomrule
\end{tabular}
\end{frame}

% Summary slide
\begin{frame}{统计与图例}
\vspace{10pt}
\centering
{\footnotesize\color{mutedgray} 1978–2024 · 以色列沃尔夫基金会颁发 \quad 交叉获奖：{\color{fieldsclr}$\blacklozenge$F} 菲尔兹奖 · {\color{abelclr}$\blacktriangle$A} 阿贝尔奖 · {\color{chernclr}$\blacksquare$C} 陈省身奖章}\\[10pt]
{\footnotesize\color{mutedgray} 首位华人得主：陈省身 (1983/84) \quad 首位女性得主：Ingrid Daubechies (2023) \quad 奖章即以陈省身命名（Chern Medal）}
\end{frame}


\end{document}
